\section{Local access and the main result}\label{sec:setting}
For a selected return branch, two endpoint laws determine a local
stationary action and hence both reflecting arcs. A global inverse still
needs to recognize repeated observations and to know when enough
branches have been observed. The completion defect solves the second
problem. The present revision addresses the first without granting
complete unknown boundaries or an already installed periodic detector.
Its additional sensor consists of certified measurements near the
contacts and flights of an actually observed short trajectory.

\subsection{The uniform class}
Write
\[
 \mathcal O=\bigcup_{i=1}^r\bigcup_{\lambda\in\Lambda}(C_i+\lambda),
 \qquad A=\covol\Lambda-\sum_{i=1}^r\area(C_i)>0.
\]
The bodies are compact and strictly convex, with disjoint closures and
positive curvature. Fix bounds $r\le r_0$, $\operatorname{diam}C_i\le D_0$,
$A_0\le A\le A_1$, $V_0\le\covol\Lambda\le V_1$, and
$\area(C_i)\ge a_0>0$. Distinct physical bodies have separation at least
$d_0>0$, and curvature has positive lower and finite upper bounds.
Some lattice presentation and its inverse are bounded, and its
fundamental parallelogram has diameter at most $b$. Neither a presentation
nor the number or identities of the body types are supplied.

After subtraction of the Steiner point, the support functions are
holomorphic and bounded by $M$ on $|\operatorname{Im}\theta|<\rho$.
For fixed $\eta,\sigma>0$ they satisfy
\begin{align}
 \|p_i(\cdot+\phi)-p_i\|_2&\ge\eta\dist(\phi,2\pi\Z),\label{eq:rot}\\
 \|p_i(-\cdot+\phi)-p_i\|_2&\ge\eta,\label{eq:refl}\\
 \inf_{O\in O(2)}\|p_i-p_{OC_j}\|_2&\ge\sigma\quad(i\ne j).\label{eq:shape}
\end{align}
Thus every body is asymmetric and different representatives are
noncongruent. In particular, $\Lambda$ is the full translation group:
a period maps a body to a congruent representative, hence to its own
lattice translate, and a compact body has no nonzero translational symmetry.

Choose $2b<R_-<R_+$. All clear shortest bridges of length at most $R_-$
are assumed to have common positive clearance and local physical margins.
Admitted additional records have length at most $R_+$ and satisfy the
same, or specified smaller, uniform margins. These include contact-chart
radius, adjacent-flight length, twist, non-grazing and time-collar bounds.
Local contact graphs extend holomorphically to a common disk with a common
bound; this also follows on smaller disks from the support and curvature
bounds. In the collar $0<d<d_*$ of a return branch, write $T=2g+d$.
There are nested common endpoint rectangles on which its action satisfies
$W-2g\le d_*/12$, and $-W_{st}\ge\omega_0>0$.
We take a closed compact class with these margins, after a narrower
analytic strip and a bounded placement gauge have been fixed. All constants
below may depend on this entire class. Such priors are substantive; they
are not conclusions of the observation.

\subsection{What the local apparatus supplies}
A speed-one experiment traces all collisions up to a bounded absolute
time. It does not initially attach obstacle-type, translate or record
labels to those collisions. Solid launch positions can be rejected. At
an encountered pair of boundary patches it has the following additional
access, called a \emph{certified local contact probe}.

It evaluates local boundary charts and their derivatives to requested
certified accuracy, measures intrinsic arclength in those charts, and
returns distance-to-solid enclosures on bounded neighborhoods of the
observed flights. It can refine these enclosures. All calls are restricted
to fixed-radius neighborhoods of the encountered contacts and of their
candidate bridge. It does not list physical bodies in an aperture, return
complete boundaries, compare complete shapes, supply period vectors, or
recognize translated copies. Section~\ref{sec:acquisition} constructs the
normal-contact and clearance test from these primitives.

For a fixed finite derivative order $K$, let $P_K(\nu)$ bound the
local measurement and refinement work for accuracy $\nu$. Finiteness
is part of the calibrated probe model; no polynomial bound is presumed.
Certified errors are deterministic bounds, not uncharged random failures.
A probabilistic probe would require an additional simultaneous error
budget. Exact computational complexity for unrestricted analytic-function
oracles is not asserted.

\begin{theorem}[Locally probed acquisition and certification]
\label{thm:main}
On the uniform class above there are an integer $K_*$ and positive
constants $\chi,p_*,e_*$ with the following properties.

\textup{(i)} A finite descriptor consisting of the gap and the contact
coefficients through degree $K_*$ identifies every oriented bounded-length
bridge up to translation \emph{among bridges of the same unknown table}.
Distinct oriented classes have descriptor distance at least $\chi$.
The order and separation bound depend only on the prior class, not on a
supplied catalogue. They do not identify an arbitrary analytic body among
all bodies in the class.

\textup{(ii)} Certified local probes of independent bounded-time free
launches build a registry of genuine records. Comparing finite descriptors
constructs each record's periodic recognition rule and its coherent sign.
No complete-body or periodic-gate installation oracle is used. An unknown
saturated witness of size at most
\begin{equation}\label{eq:witness-main}
 w_0=r_0+1+\lfloor\log_2 Q\rfloor,\qquad
 Q=\max\{1,\lceil B_0^2/V_0\rceil\},\quad
 B_0=(2r_0-1)(R_++2D_0),
\end{equation}
is discovered with per-record conditional hazard at least $p_*>0$.

\textup{(iii)} For every $\xi>0$ and $0<\delta<1/4$, two growing
endpoint histograms per registered record give simultaneous whole-table
certification at all epochs with probability at least $1-\delta$.
For the first certificate of accuracy $\xi$, denoted $T_\xi$, and all
sufficiently large deterministic $k\ge1$,
\begin{equation}\label{eq:main-tail}
 \Prob\{T_\xi>k\}\le\frac{\delta}{8(k+1)^a}+w_0e^{-p_*k},
 \qquad a=6.
\end{equation}
Consequently $T_\xi$ is almost surely finite and the expected number
of launch attempts is finite. The reconstruction error is measured in
support-function $C^2$ and a period basis, with one common isometry,
permutation and choice of representatives and lattice basis.

Through epoch $m$ there are $O(m^4)$ accepted preparations. At most
$O(m^4+m\log(m/\delta))$ launch attempts suffice with the stated
high probability. The largest square half-width is
$O(m^{(\beta+4)/(\beta+3)})$, where $0<\beta\le1$ is the density
regularity exponent. Local probe work is counted separately as
\eqref{eq:probe-cost}. If $P_{K_*}(\nu)\le C\nu^{-s}$, it is
$O(m^{4+s})$; choosing $a>4+s$ makes its expectation finite as well.
Analytic fitting, apparatus travel and setup are not included in the
launch count. No end-to-end polynomial algorithm is claimed.
\end{theorem}

The periodic rule in part (ii) is an output of local identification,
not a hidden input. Nevertheless the theorem remains local-probe-relative.
A passive unmarked trajectory without boundary or clearance measurements
is a different observation model. The earlier geometric-oracle theorem
continues to apply under its original hypotheses and is retained in full.

\subsection{Records and normalization}
For one clear normal bridge select a three-impact return to its source.
At two strictly active absolute times the ideal record consists of two
subprobability densities $f_1(s,t),f_2(s,t)$ of the first and last
intrinsic arclengths from the source normal contact. Both use one common
sign, with simultaneous reversal allowed. Failure and endpoints outside
the measured rectangle remain outcomes. Two windows mean two functions,
or two increasingly fine two-dimensional histograms, not two numbers.
The target normalization is $\dd q\dd\vartheta/(2\pi A)$.
Laboratory-square sampling approximates this without a known cell.
Neither conditioning on return success nor normalizing a local aperture
would preserve the unseen-area information used by the certificate.
